% ЗАКЛЮЧЕНИЕ, без номер. Опира се само на измереното.
\section*{ЗАКЛЮЧЕНИЕ}
\label{sec:conclusion}
\addcontentsline{toc}{section}{ЗАКЛЮЧЕНИЕ}

Изградени са демон за маршрутизация по състояние на връзките и дискретно-събитиен симулатор,
който изпълнява същата протоколна логика върху управлявано време, така че отказът на връзка да
бъде възпроизводимо събитие. Протоколното ядро няма достъп до часовника и до сокетите, а работи
над абстракция от пет операции, така че едни и същи изходни файлове се изпълняват и в симулация,
и над реални интерфейси, и измереното поведение е свойство на реализацията.

\textbf{Резултатите} са четири твърдения, всяко от 250 валидни измервания и нула отхвърлени.
Първо, времето за сходимост се определя почти изцяло от откриването на отказа: следва интервала
за отпадане с отношение от 0,81 до 0,98 и нараства само с 2,0 ms на хоп, деветнадесет пъти
по-малко от разсейването между повторенията. Второ, при преминаване от пръстен към пълна мрежа
с еднакъв брой възли връзките растат 5,5 пъти, съобщенията 10,8 пъти, а байтовете 35,2 пъти,
докато промените в маршрутните таблици намаляват от 42 на 2. Трето, преходните цикли зависят
преди всичко от топологията, а признакът е дали алтернативният път обръща оцеляло ребро.
Четвърто, едностранният отказ се възстановява от 153 до 206 ms по-бавно от двупосочния.

\textbf{Предимствата} спрямо измереното са три: пълна възпроизводимост, потвърдена от 500
изпълнения с идентична контролна сума; достатъчност на абстракцията, тъй като същият клас
\code{Router} работи и над виртуален часовник, и над реални UDP сокети; и липса на външни
зависимости. \textbf{Недостатъците} са изброените в раздел~\ref{sec:results}, плюс
нереализираното записване на маршрутите в операционната система.

Два дефекта заслужават отбелязване, защото и двата бяха намерени от измерване, а не от
разсъждение. Първият беше запазен обзор на базата, който при възпроизвеждане описваше по-ранно
състояние и оставяше една връзка невидима за цялата мрежа. Вторият беше заявка за обявление,
чийто отговор пристига по-стар от вече наученото отдругаде и затова не изчистваше чакането,
оставяйки съседството в междинно състояние завинаги. И двата се проявяваха само при определени
зърна и щяха да останат незабелязани при неповторимо изпълнение, което е практическото
потвърждение на изходната теза. Посоки за по-нататъшна работа са кодиране по формата на RFC
2328, което би позволило сверяване с FRRouting или BIRD, записване на маршрутите в
операционната система, модел на канал със загуби и по-бързо откриване на
отказ~\cite{rfc5880}.
